\section{Bốn lần thử trước ViT}
\label{sec:ch10-truoc-vit}

\index{thiên kiến quy nạp!bốn hướng thử trước ViT}%
Từ 2018 tới 2020, giữa bài Transformer và bài ViT, có bốn cách người ta đem
attention vào thị giác. Đọc chúng cạnh nhau thì thấy chúng không phải bốn ý
tưởng rời rạc mà là bốn vị trí trên cùng một trục: giữ lại bao nhiêu phần của
hai ràng buộc ở mục~\ref{sec:ch10-hai-rang-buoc}.

Và cả bốn đều vấp vào cùng một bức tường trước, bức tường của
chương~\ref{ch:chi-phi-va-song-song}: chi phí bậc hai theo số vị trí. Với văn
bản, một câu là vài chục token. Với ảnh, số vị trí là số pixel.

\begin{bridge}{Image Transformer, Parmar và cộng sự 2018}
Cách thứ nhất: giữ attention, nhưng \emph{bắt} nó cục bộ.

Bài nói thẳng lý do ở mục 3.3, và con số trong câu ấy là chỗ chương này gặp lại
chương~\ref{ch:chi-phi-va-song-song}~\autocite{parmar2018image}:

\begin{quote}
\enquote{The decoders in our experiments, however, produce 32x32 pixel images
with 3072 positions, rendering attending to all positions impractical.}
\end{quote}

3072 là \(32 \times 32 \times 3\). Một ảnh nhỏ xíu, và attention toàn cục trên
nó đã không chạy nổi.

Nên bài chia ảnh thành khối truy vấn, mỗi khối kèm một khối bộ nhớ lớn hơn bao
quanh, và mỗi vị trí chỉ nhìn trong khối bộ nhớ của mình. Đó chính là ràng buộc
\tn{tính cục bộ}{locality}, đặt lại bằng tay lên một cơ chế vừa mới bỏ nó đi.

Kết quả: log-likelihood âm trên ImageNet từ 3.83 xuống 3.77, và trên CIFAR-10
là 2.90 so với 2.85 của PixelSNAIL. Tức là hơn ở một bộ dữ liệu và kém ở bộ
kia.
\end{bridge}

Cách thứ hai đi xa hơn: bỏ hẳn \tn{tích chập}{convolution}, thay bằng
self-attention cục bộ ở mọi tầng. Ramachandran và cộng sự làm đúng thế trên ResNet, và bảng 1 của họ đáng
nhìn kỹ~\autocite{ramachandran2019standalone}. Ở ResNet-50, mô hình chỉ-attention
đạt 77.6\% so với 76.9\% của baseline, với 18.0 triệu tham số so với 25.6
triệu, và 7.2 tỷ phép tính dấu chấm động - floating point operation, viết tắt
là FLOP - so với 8.2 tỷ.

Hơn baseline, và rẻ hơn. Đọc nhanh thì đây là bằng chứng attention thắng tích
chập. Đọc kỹ thì cái họ thay vào vẫn cục bộ: mỗi vị trí vẫn chỉ nhìn một lân
cận. Cái bị bỏ là \tn{chia sẻ trọng số}{weight sharing} theo kiểu kernel cố
định, còn ràng buộc thứ nhất thì giữ nguyên. Nên bảng ấy nói rằng \emph{tính cục bộ} là phần đáng giữ,
chứ không nói rằng hai ràng buộc đều bỏ được.

Cách thứ ba là thỏa hiệp: giữ tích chập, thêm attention vào cạnh nó. Bello và
cộng sự báo cáo hơn 1.3\% top-1 trên ImageNet so với ResNet50, và 1.4 điểm
mean average precision, viết tắt là mAP, trên
COCO~\autocite{bello2019attention}. Không có gì bị bỏ đi ở đây, nên nó không
trả lời câu hỏi của chương này; nó có mặt vì nó là thứ phần lớn người làm thực
tế đã dùng trong hai năm ấy.

\begin{bridge}{Image GPT, Chen và cộng sự 2020}
Cách thứ tư: không thỏa hiệp gì cả. Lấy nguyên GPT của
chương~\ref{ch:tien-huan-luyen}, cho nó ăn pixel theo thứ tự raster như thể ảnh
là một chuỗi token, và trả giá~\autocite{chen2020igpt}.

Không có tính cục bộ. Không có chia sẻ trọng số theo vị trí. Không có bất cứ
thứ gì trong ảnh được nói cho mô hình biết.

Nó chạy được. Trên CIFAR-10, linear probe đạt 96.3\% và fine-tune đạt 99.0\%;
trên ImageNet, linear probe đạt 72.0\%.

Cái giá thì bài nói ở hai chỗ. Mục Introduction ghi phần cứng: 2048 lõi tensor
processing unit, viết tắt là TPU. Và
mục Discussion tự đánh giá:

\begin{quote}
\enquote{we observed that our approach requires large models in order to learn
high quality representations. iGPT-L has 2 to 3 times as many parameters as
similarly performing models on ImageNet and uses more compute.}
\end{quote}

Bốn mô hình của bài lần lượt 76 triệu, 455 triệu, 1362 triệu và 6801 triệu tham
số.

Tôi để bài tự nói câu ấy chứ không đi tìm một con số giờ-GPU ở nơi khác: con số
hay được trích lại cho bài này nằm trên một trang blog mà tôi không mở được, và
một nguồn tôi không đọc được thì không vào sách.
\end{bridge}

\subsection*{Cái bốn hướng ấy nói chung một câu}

\index{thiên kiến quy nạp!kết luận chung của bốn hướng}%
Xếp bốn hướng theo lượng \tn{thiên kiến quy nạp}{inductive bias} còn giữ lại thì ra
một dãy: Image
Transformer giữ tính cục bộ và bỏ kernel cố định; Stand-Alone giữ tính cục bộ
và thay cách trộn; Attention Augmented giữ cả hai và thêm; iGPT bỏ hết.

Và kết quả đi theo dãy ấy đúng như bạn đoán, với một chỗ ngoặt: iGPT, cái bỏ
hết, \emph{vẫn chạy được}, chỉ là đắt tới mức lố bịch.

Đó là tiền đề mà chương~\ref{ch:vit} cần và
chương~\ref{ch:tien-huan-luyen} đã dựng một nửa. Thiên kiến quy nạp không phải thứ
bắt buộc phải có. Nó là thứ mua được bằng dữ liệu và bằng tính toán, và câu hỏi
duy nhất là tỷ giá.

Mục sau là chỗ tôi đo tỷ giá ấy, ở quy mô một cái laptop.
